\documentclass[10pt,a4paper]{amsart}
\usepackage[margin=19mm]{geometry}
\usepackage{amsmath,amssymb,mathtools}
\usepackage[scr=rsfs,cal=euler]{mathalfa}
\usepackage{xfrac}
\usepackage[unicode,hidelinks,bookmarks=false]{hyperref}
\newcommand{\ie}{i.e.\ }
\newcommand{\cf}{cf.\ }
\newcommand{\resp}{respectively\ }
\newcommand{\Subsection}{\subsection}
\providecommand{\nobreakdash}{\nobreak\hbox{-}}
\setlength{\emergencystretch}{2em}
\multlinegap0pt
\usepackage{amssymb}
\usepackage{tikz}
\newtheorem{prop}{Proposition}
\newtheorem{cl}{Claim}
\def \N {\mathbb N}
\newcommand{\Nzero}{\mathbb N_0}
\title{An affirmative answer to Owings's sumset question}
\author{Wen Huang, Zhengxing Lian, Song Shao, Rongzhong Xiao, Leiye Xu,, Shuhao Zhang}
\date{Source: arXiv:2607.17333v3 (CC BY 4.0)}
\begin{document}
\maketitle

\noindent\emph{Source and licence.} An affirmative answer to Owings's sumset question. Version arXiv:2607.17333v3 (CC BY 4.0), distributed by arXiv under the Creative Commons Attribution 4.0 International licence, http://creativecommons.org/licenses/by/4.0/, as recorded in the arXiv licence metadata for this exact version and shown on the abstract page https://arxiv.org/abs/2607.17333v3. Full licence notice: \texttt{licenses/arxiv-2607.17333-CC-BY-4.0.txt}.\par\smallskip

\noindent\textit{Editorial scope.} Counted unit: the article's prop prop6-2 (source0.tex lines 2893--2903). The article's complete original proof of it is delivered with it (source0.tex lines 2904--3040, one \texttt{proof} environment of 137 lines). No mathematics is written, repaired or re-derived: the statement and the proof are the article's own lines, in the article's own notation, and no retained line is substituted or re-flowed.  No further source-local context is needed: every cross-reference of every retained line resolves inside the two delivered spans.  Nothing was omitted from any retained span, so the package carries no omission record.  No retained line contains a citation, so the wrapper carries no bibliography and no bibliography entry.  No retained line contains a figure environment, an image-inclusion command or drawing code, so no image is delivered and the package declares no asset dependency.  Editorial formatting only, in addition: an AMS \texttt{amsart} preamble that restates the article's own declarations and macros used by the retained lines, namely the environments \texttt{prop}, \texttt{cl} on separate counters, together with the standard packages those lines need.  The retained environments print in this document as Proposition 1, Cl 1 in order of appearance, whereas the article prints the same items with the numbering of its own full document.  The complete native source of the article as distributed in the arXiv e-print for this exact version is bundled unchanged as \texttt{sources/205559/source0.tex}.
\begin{prop}\label{prop6-2}
For any $\ell,m\in \N$ with $\ell\neq m$, there is a $2$-coloring of
$\N$ such that, for every infinite $B\subset\N$ and all
$t_1,t_2\in\mathbb Z$, the set
\[
\{mx+\ell y+t_1:x,y\in B,\ x<y\}
\cup
\{mx+\ell y+t_2:x,y\in B,\ x>y\},
\]
whenever contained in $\N$, is not monochromatic.
\end{prop}

\begin{proof}
Exchanging $m$ and $\ell$ and simultaneously interchanging the two ordered
pieces if necessary, it suffices to consider $\ell<m$. Put
$r=m/\ell>1$ and color $(0,\infty)$ as follows:
\begin{itemize}
	\item color $1$: $[r^{2n},r^{2n+1})$ for all $n\in\Nzero$;
	\item color $2$: $[r^{2n+1},r^{2(n+1)})$ for all $n\in\Nzero$;
	\item color $1$: otherwise (that is, $(0,1)$).
\end{itemize}
Restrict this coloring to $\N$ and denote it by $\phi$.

\medskip
Assume that there exist an infinite $B\subseteq \N$ and
$t_1,t_2\in\mathbb Z$ such that $\phi$ takes a constant value
$c\in\{1,2\}$ on
\[
\{mx+\ell y+t_1:x,y\in B,x<y\}
\cup\{\ell x+my+t_2:x,y\in B,x<y\}.
\]
Here the second family is the $x>y$ family in the statement after the two
variables have been interchanged.

We choose $n_1,n_2,k$ successively. Since
\[
 rm-\ell=\frac{m^2-\ell^2}{\ell}>0,
\]
we may first choose $n_1\in B$ so large that $mn_1+t_1>0$ and
\begin{equation}\label{eq:asym-n1}
 \ell n_1+t_2<r(mn_1+t_1).
\end{equation}
Next choose $n_2\in B$, $n_2>n_1$, so large that
\begin{equation}\label{eq:asym-n2}
 r(mn_1+t_1)<\ell n_2+t_2.
\end{equation}
Finally, choose $k\in B$, $k>n_2$, sufficiently large. For this choice
the following conditions hold:
\begin{itemize}
\item[(1)] $\ell n_2+mk+t_2,\ell n_1+mk+t_2,\ell k+mn_1+t_1>1$.
\item[(2)] $r(mn_{1}+t_1)<\ell n_2+t_2$.
\item[(3)] $1<(\ell n_2+mk+t_2)/(\ell n_1+mk+t_2)<r,1<(\ell n_1+mk+t_2)/(\ell k+mn_1+t_1)<r$.
\end{itemize}
Indeed, (1) holds for all sufficiently large $k$, while (2) is
\eqref{eq:asym-n2}. Moreover,
\[
 \frac{\ell n_2+mk+t_2}{\ell n_1+mk+t_2}\longrightarrow1
 \qquad\text{as }k\to\infty.
\]
The quotient is greater than $1$ because $n_2>n_1$, and hence it is
less than $r$ for all sufficiently large $k$. Similarly,
\[
 \frac{\ell n_1+mk+t_2}{\ell k+mn_1+t_1}\longrightarrow
 \frac m\ell=r.
\]
The quotient is less than $r$ by \eqref{eq:asym-n1}, and it is greater
than $1$ for all sufficiently large $k$, since $m>\ell$. This proves
(3), and the required $k$ exists because $B$ is unbounded.
Next, we introduce a claim.
\begin{cl}
	If $x,y>1$, $1<x/y<r$, and $\phi(x)=\phi(y)$, then $x,y\in [r^d,r^{d+1})$ for some $d\in \Nzero$.
\end{cl}
\begin{proof}
	Suppose not. At this point, we take $d\in \Nzero$ such that $x\in [r^d,r^{d+1})$ and $y\notin [r^d,r^{d+1})$. By $1<x/y<r$, $r^{d-1}\le y<r^{d}$. This contradicts $\phi(x)=\phi(y)$. Therefore, $x,y\in [r^d,r^{d+1})$ for some $d\in \Nzero$.
\end{proof}
Now, assume that $(\ell n_1+mk+t_2)\in [r^d,r^{d+1})$ for some $d\in \Nzero$. By the above claim, $(\ell n_2+mk+t_2),(\ell k+mn_1+t_1)\in [r^d,r^{d+1})$. Then $r^{d+1}>(\ell n_2+mk+t_2)$ and $(\ell k+mn_1+t_1)r\ge r^{d+1}$. This contradicts (2). This completes the proof.
%Choose $n_1\in B$ so large that $|t_1|+|t_2|+1\le n_1$. Since $B$ is
%unbounded, we may then choose $n_2\in B$, $n_2>n_1$, so large that
%\begin{equation}\label{sq-est-1}
%m^2n_1+mt_1<\ell^2n_2+\ell t_2
%\quad\text{and}\quad
%\ell n_1+t_2<mn_2+t_1.
%\end{equation}
%Then
%\[
%0<\min\{mn_1+t_1,\ell n_1+t_2,mn_2+t_1,\ell n_2+t_2\}.
%\]
%Next choose $k\in B$ sufficiently large that $k>n_2$ and, for the unique
%$d\in\Nzero$ satisfying
%\begin{equation}\label{lk-est-1}
%r^d\le \ell k<r^{d+1},
%\end{equation}
%we have
%\begin{equation}\label{n1n2-est-1}
%(m+\ell)(n_1+n_2)+|t_1|+|t_2|<r^d(r-1).
%\end{equation}
%Such a choice is possible because $B$ is unbounded and the corresponding
%exponent $d$ tends to infinity with $k$.
%
%By  \eqref{n1n2-est-1} and $ r>1$ we have
%$0<mn_1+t_1<mn_2+t_1<r^{d+2}-r^{d+1}$. By \eqref{lk-est-1} we have
%\begin{equation}\label{lk-est-2}
%r^d<mn_1+\ell k+t_1<mn_2+\ell k+t_1<r^{d+2}.
%\end{equation}
%By \eqref{n1n2-est-1} we have 
%$0<\ell n_1+t_2<\ell n_2+t_2<r^{d+3}-r^{d+2}$. By \eqref{lk-est-1} we have
%$r^{d+1}\le mk=r\ell k<r^{d+2}$ and
%\begin{equation}\label{mk-est-1}
%r^{d+1}<\ell n_1+mk+t_2<\ell n_2+mk+t_2<r^{d+3}.
%\end{equation}
%
%There are two cases. First suppose
%$mn_2+\ell k+t_1\in(r^d,r^{d+1})$. Using the equality of colors and
%\eqref{mk-est-1}, we obtain
%\[
%\ell n_1+mk+t_2\in[r^{d+2},r^{d+3}).
%\]
%Combining this with $mn_2+\ell k+t_1<r^{d+1}$ gives
%\begin{equation}\label{con-est-1}
%(\ell n_1+mk+t_2)-(mn_2+\ell k+t_1)>r^{d+2}-r^{d+1}=r^{d+1}(r-1).
%\end{equation}
%On the other hand,
% \begin{align*}
%(\ell n_1+mk+t_2)-(mn_2+\ell k+t_1)
%&=(m-\ell)k+\ell n_1+t_2-(mn_2+t_1)<(m-\ell)k\\
%&=(r-1)\ell k<r^{d+1}(r-1),
%\end{align*}
%which contradicts \eqref{con-est-1}.
%
%Now suppose $mn_2+\ell k+t_1\in[r^{d+1},r^{d+2})$. Using the equality
%of colors together with \eqref{lk-est-2} and \eqref{mk-est-1}, we obtain
%\[
%mn_1+\ell k+t_1\in[r^{d+1},r^{d+2})
%\quad\text{and}\quad
%\ell n_2+mk+t_2\in[r^{d+1},r^{d+2}).
%\]
%Therefore
%\[
%r\bigl(r^{d+1}-(mn_1+t_1)\bigr)
%\le r\ell k=mk
%<r^{d+2}-(\ell n_2+t_2).
%\]
%After multiplying by $\ell$ and simplifying, this gives
%$\ell^2n_2+\ell t_2<m^2n_1+mt_1$, contrary to \eqref{sq-est-1}.
%
%Both cases lead to a contradiction. Therefore the displayed configuration,
%whenever contained in $\N$, is not monochromatic for any infinite
%$B\subseteq\N$ and any $t_1,t_2\in\mathbb Z$.
\end{proof}

\end{document}
